%
% Tesi D.S.I. - modello preso da
% Stanford University PhD thesis style -- modifications to the report style
%
%%%%%%%%%%%%%%%%%%%%%%%%%%%%%%%%%%%%%%%%%%%%%%%%%%%%%%%%%%%%%%%%%%%%%%%%%%%
%                                                                         %
%			TESI DOTTORATO                                                   %
%			______________                                                   %
%                                                                         %
%			AUTORE: Elena Pagani                                             %
%                                                                         %
%			Ultima revisione: 7.X.1998                                       %
%           correzioni atrent                                             %
%%%%%%%%%%%%%%%%%%%%%%%%%%%%%%%%%%%%%%%%%%%%%%%%%%%%%%%%%%%%%%%%%%%%%%%%%%%
%
%
\documentclass[a4paper,12pt]{report}
%    \renewcommand{\baselinestretch}{1.6}      % interline spacing
%
% \includeonly{}
%
%			PREAMBOLO
%
\usepackage[a4paper]{geometry}
\usepackage{amssymb,amsmath,amsthm}
\usepackage{graphicx}
\usepackage{url}
\usepackage{epsfig}
\usepackage[english]{babel}
\usepackage{setspace}
\usepackage{tesi}
\usepackage{algorithm}
\usepackage{algcompatible}
%\usepackage{algorithmicx}
\usepackage{comment}
\usepackage[section]{placeins}
\usepackage[a-1b]{pdfx}
\usepackage[pdfa]{hyperref}
\usepackage[utf8]{inputenc}
\usepackage{csquotes}
\usepackage{tikz}
\usepackage{subcaption}
\usepackage{tabularx}
\usepackage{longtable}
\usepackage{booktabs}
\usepackage{mathtools}
%\usepackage{float}
%\restylefloat{table}
\newif\ifshowtableLinear
\showtableLinearfalse % Cambia in \showtabletrue per mostrare la tabella

\setcounter{secnumdepth}{5}

%
\newtheorem{myteor}{Teorema}[section]
%
\newenvironment{teor}{\begin{myteor}\sl}{\end{myteor}}
%
\usepackage[
    backend=biber,
    style=authoryear-comp,
    sorting=ynt,
    maxcitenames=6,
    maxbibnames=10
]{biblatex}
\addbibresource{biblio.bib}

\usepackage{subfiles}
%
%			TITOLO
%
\begin{document}

\includegraphics[width=\textwidth]{thesis/picture/tesiSCIENZE_TECNOLOGIE.jpg}
\title{Methods for traffic management in urban air mobility}
\author{Filippo Magi}
\dept{Corso di Laurea in Informatica Magistrale} 
\anno{2023-2024}
\matricola{17305A}
\relatore{Prof. Alberto Ceselli}
\correlatore{Prof. Claudia D'Ambrosio}
%
%        \submitdate{month year in which submitted to GPO}
%		- date LaTeX'd if omitted
%	\copyrightyear{year degree conferred (next year if submitted in Dec.)}
%		- year LaTeX'd (or next year, in December) if omitted
%	\copyrighttrue or \copyrightfalse
%		- produce or don't produce a copyright page (false by default)
%	\figurespagetrue or \figurespagefalse
%		- produce or don't produce a List of Figures page
%		  (false by default)
%	\tablespagetrue or \tablespagefalse
%		- produce or don't produce a List of Tables page
%		  (false by default)
% 
%			DEDICA
%
\beforepreface
\prefacesection{}
        {\hfill \Large {\sl dedicato a tutti coloro che mi sono stati vicino nei momenti difficili}}
% 
%			PREFAZIONE
%
\prefacesection{Preface}
\subfile{thesis/sections/intro/preface}
%
%
%			ORGANIZZAZIONE
\section*{Thesis Organization}
\label{organizzazione}

This thesis is organized as follows:
\begin{itemize}
    \item \textbf{Chapter 1}: We introduce the problem and review previous works.
    \item \textbf{Chapter 2}: We present a general-use Mixed integer linear model to define the problem, highlight the differences with previous works.
    \item \textbf{Chapter 3}: We introduce a new model for solving tactical deconfliction with some fixed paths.
    \item \textbf{Chapter 4}: We present a heuristic approach for solving tactical deconfliction while allowing path modifications.
    \item \textbf{Chapter 5}:We describe the instances used for testing and explain how they were generated.
    \item \textbf{Chapter 6}: We analyze the computational results obtained by both the mathematical model and the heuristic for solving tactical deconfliction.
\end{itemize}
% 
% 

\afterpreface
%			CAPITOLO 1: dshjkfg
\chapter{Introduction}
\subfile{thesis/sections/intro/Introduction}
\subfile{thesis/sections/intro/pelegrinSummary}
\subfile{thesis/sections/intro/PortoleauSummary}

%
%
\chapter{Definition of the problem}
%\newpage
\section{The model}\label{sec:modelDescription}
We define the problem using its mathematical model.  
First, Section \ref{sec:modelData} explain the input data, followed by Section \ref{sec:modelVar} which introduces the decision variables.
%,Section \ref{sec:modelData}, which explains the input data.  
Next, we present the constraints in Section \ref{ssec:bigModel:constrains}, and finally, in Section \ref{sec:linearization}, we describe the techniques used to linearize the nonlinear constraints.  
\subfile{thesis/sections/generalModel/variablesAndData}
\subfile{thesis/sections/generalModel/constraints}
\subfile{thesis/sections/generalModel/linearization}
\subfile{thesis/sections/mercedesChanges}
%\section{Testing maps}\label{sec:maps}

\chapter{Exact Method}\label{chap:ExactAlgo}
\subfile{thesis/sections/PelegrinModification}

\chapter{Heuristic Algorithms}\label{chap:HeuristicAlgo}
\subfile{thesis/sections/heuristic}

\chapter{Maps used and instance creation}\label{sec:maps}
\subfile{thesis/sections/maps/topologies}
\subfile{thesis/sections/maps/generationOfIstances}

\chapter{Computational results}\label{chap:results}
In this chapter, we present the results obtained for both the model and the heuristic.  
For a complete overview of all instance results, refer to \cite{Magi_Methods_for_Traffic_2025}.  
In both cases, we performed our calculations on the Atlas server at Polytechnique, equipped with an Intel\textregistered~Xeon\textregistered~Platinum 8362 CPU @ 2.80GHz and 2TiB of system memory.
\subfile{thesis/sections/conclusions/PelegrinConclusion}
\subfile{thesis/sections/conclusions/HeuristicConclusion}
\chapter{Conclusion}\label{chap:conclusion}
\subfile{thesis/sections/conclusions/Conclusion}
%
%			BIBLIOGRAFIA

\printbibliography

\appendix 
\chapter{Models} 
\section{The non linearized model}\label{sec:wholeModel}
\subfile{thesis/models/wholeModel} 
\section{The linearized model}\label{sec:linearModel}
\subfile{thesis/models/linearModel}
\section{Pelegrin Model with Path Changing Capability}\label{sec:MercedesModifedModel}
\subfile{thesis/models/MercedesModel}
\chapter{The algorithms}
\section{The matheuristic}\label{alg:matheuristic}
\subfile{thesis/pseudocode/matheuristic}
\chapter{Tables}
\section{Pelegrin's instances}
%number of variables
\subsection{Number of variables for fixed flights}
\subfile{thesis/tables/mercedes/variables/SummaryTablesFixed}
\subsection{Number of variables for free flights}
\subfile{thesis/tables/mercedes/variables/SummaryTablesFree}
\subsection{Comparison between number of variables for free and fixed flights}
\subfile{thesis/tables/mercedes/variables/ComparisonTables}
%times and results
\subsection{Results and calculation time for fixed flights}
\subfile{thesis/tables/mercedes/times/SummaryTablesFixed}
\subsection{Results and calculation time for free flights}
\subfile{thesis/tables/mercedes/times/SummaryTablesFree}
\subsection{Comparison between results and calculation time for free and fixed flights}
\subfile{thesis/tables/mercedes/times/ComparisonTablesResults}
\section{Results for the matheuristic}
\subsection{Summary tables}
\subfile{thesis/tables/heuristic/times/SummaryTables}\label{chap:heuristic:tables}
\subsection{Comparison between heuristic and no fixed path}
\subfile{thesis/tables/heuristic/times/comparison}
\end{document}


 
